\subsection{Degeneracies with Standard Astrophysical Effects}
  \label{subsec:degeneracies}

  Several mechanisms within the $\Lambda$CDM framework
  have been proposed to address galaxy rotation curves and the Hubble tension.
  It is therefore necessary to examine possible degeneracies.

  At galactic scales, feedback processes and baryon--halo coupling can reproduce
  some features of flattened rotation curves within dark matter halo models,
  particularly in low-mass systems~\cite{DiCintio2014Feedback}.
  Such mechanisms typically require galaxy-dependent tuning and do not
  naturally reproduce the observed correlation between surface density and
  dynamical response across the full range of disk galaxies~\cite{McGaugh2016}.

  In contrast, the effective saturation framework introduces no halo degrees of
  freedom and relies on a single universal scale.
  All diversity in rotation curve shapes arises from the observed baryonic distributions.
  This leads to distinct residual patterns and scaling relations,
  which can be used to discriminate between models in detailed rotation curve analyses.

  At cosmological scales, local inhomogeneities, peculiar velocities, and sample
  variance are known to affect Hubble constant measurements~\cite{Freedman2021TRGB}.
  While these effects contribute to statistical scatter,
  multiple studies indicate that they are generally insufficient
  to account for the full observed tension~\cite{Verde2019Review}.
  The effective saturation framework instead predicts a systematic environment-dependent bias, reflecting a change in
  effective transition resolution rather than a stochastic fluctuation.

  A key aspect of the present framework is that the fundamental bound $b$ is universal.
  It fixes a constant maximal resolvability of the effective projection,
  which may be viewed as a uniform ``pixel density'' of the geometric description of the Universe.

  Importantly, different environments do not exploit this fixed resolution in the same way.
  Dense systems, such as galaxies or clusters, concentrate many effective
  relations within a limited physical volume.
  In these regimes, the available resolution is primarily consumed by internal
  structure, and the effective transition rate remains close to the Newtonian or background behavior.

  By contrast, cosmic voids distribute the same fixed resolution over much larger volumes.
  The effective density of contributing relations is therefore reduced, and the
  projection operates closer to its saturation threshold.
  This leads to an enhanced effective coupling in low-density environments,
  without any change in the underlying bound $b$ or in the global expansion history.

  This distinction cannot be absorbed into standard astrophysical degeneracies.
  It reflects a difference in how environments of contrasting density sample
  a fixed geometric resolution, rather than a reparameterization of existing dynamical models.
